\subsection*{All four-point tests and a nontrivial common kernel}

\statement{B.18}{Proved}{An exact list of the 495 cross-ratio pairs}
Let \(T\leq\mathrm{PGL}_2(\mathbb Q(i))\) be generated by
\[
 u\longmapsto-u,\qquad u\longmapsto1/u,\qquad
 u\longmapsto\frac{u+i}{u-i}.
\]
It is tetrahedral of order twelve. Label its twelve elements
in the order of the breadth-first construction in
\texttt{code/crossratio\_495.py}, and write
\(g_j(u)=(a_ju+b_j)/(c_ju+d_j)\).
The ordered tetrahedral parameter curve is
\[
 S_T=\mathbb P^1\setminus E,\qquad
 E=\{\infty\}\cup
 V\bigl(u(u^4-1)(u^8+14u^4+1)\bigr);
\]
it has fourteen punctures and fundamental group \(F_{13}\).
There is an explicit labelling of
\[
 S_5=\mathbb P^1\setminus\bigl(\{0,\infty\}\cup\mu_5\bigr),
 \qquad (0,\infty,1,\omega,\ldots,\omega^4,
                       t,t\omega,\ldots,t\omega^4)
\]
which agrees at an icosahedral basepoint with the tetrahedral labelling.
For each increasing four-subset \(q=(i,j,k,l)\), the attached
\texttt{data/crossratio\_495.csv} gives the rational functions
\(\psi_q:S_T\to M_{0,4}\) and \(\phi_q:S_5\to M_{0,4}\), with that
common basepoint. We use the ordered covers: on the unlabelled
quotient tripods a fixed four-subset need not be invariant under
the label monodromy.
\begin{proof}
The listed fractional transformations have precisely twelve products
up to scalar. Their exceptional orbits have sizes six, four and four.
The six-point orbit is \(\{0,\infty,\pm1,\pm i\}\); the remaining
eight points are the roots of \(u^8+14u^4+1\).
The program verifies the group closure exactly.

For homogeneous points \(P_j=(A_j:B_j)\), put
\[
 [P_i,P_j;P_k,P_l]
   =\frac{\det(P_i,P_k)\det(P_j,P_l)}
           {\det(P_i,P_l)\det(P_j,P_k)}.
\]
For the tetrahedral curve take
\(P_j=(a_ju+b_j:c_ju+d_j)\); for the cyclic curve use the
homogeneous points listed in
\texttt{data/crossratio\_labels.csv}. These formulas define all
495 pairs without path ambiguity at the level of pointed curves.
No four points collide on either parameter curve, so the images
avoid \(0,1,\infty\).

For the basepoint check work in
\[
 \mathbb Q(q)/(q^8-q^6+q^4-q^2+1),\qquad q=e^{\pi i/10},
 \quad i=q^5,\quad\omega=q^4.
\]
Take \(u_0=q^7+q^6+q^4+q^3\).
Projectively normalize the first, third and second tetrahedral
points to \(0,\infty,1\), respectively. The two five-element
orbits then have the displayed form with
\(t_0=q^6-q^4-2\). The program verifies all 495 equalities
\(\psi_q(u_0)=\phi_q(t_0)\) by exact arithmetic in this
cyclotomic field. It finds 351 tetrahedral maps of degree four
and 144 of degree two. The configuration has both tetrahedral
and fivefold symmetry; its finite subgroup of
\(\mathrm{PGL}_2(\mathbb C)\) is therefore icosahedral.
\end{proof}

\statement{B.19}{Proved}{The paired four-point equations have nontrivial discrete solutions}
There is an explicit nonidentity word \(w\in\pi_1(S_T,u_0)\)
such that
\[
  \psi_{q*}(w)=1=\phi_{q*}(1)
       \qquad\text{for all }q\in\binom{\{1,\ldots,12\}}4.
\]
The same equalities hold after full profinite completion.
Thus these equations alone do not characterize the icosahedral
centralizer: they already admit noncentralizing discrete
tetrahedral braids.
\begin{proof}
Each determinant in the cross-ratio formula is a polynomial of
degree at most two. Hence \(\deg\psi_q\leq4\), independently of
the machine count. A degree-\(d\) rational function has at most
\(3d\) distinct points over \(0,1,\infty\).
Since \(E\) has fourteen points, \(\psi_q\) extends across at
least two of the punctures of \(S_T\). The induced homomorphism
on fundamental groups therefore kills a meridian about at least
one of these fourteen punctures.

Choose positively oriented, whiskered meridians
\(e_1,\ldots,e_{14}\) with
\(e_1\cdots e_{14}=1\). For a list of words define recursively
\[
\begin{split}
 W(v_1)&=v_1,\\
 W(v_1,\ldots,v_m)
  &=[W(v_1,\ldots,v_{\lfloor m/2\rfloor}),
     W(v_{\lfloor m/2\rfloor+1},\ldots,v_m)],
\end{split}
\]
where \([a,b]=aba^{-1}b^{-1}\).
Put \(w=W(e_1,\ldots,e_{14})\).
Killing any one argument kills its subtree and then the entire
word. Consequently every \(\psi_{q*}\) kills \(w\).

This word is nontrivial. Substituting
\(e_{14}=(e_1\cdots e_{13})^{-1}\) leaves a freely reduced word
of length 384, included in
\texttt{data/crossratio\_common\_kernel\_word.txt}.
For an independent finite witness,
\texttt{code/check\_crossratio\_kernel.py} lists thirteen
matrices in \(\mathrm{SL}_2(\mathbb F_{101})\); evaluation gives
\[
 w\longmapsto
 \begin{pmatrix}5&32\\72&57\end{pmatrix}\ne I .
\]
The determinant is \(1\) modulo \(101\).
Thus the word remains nontrivial in the full profinite completion.

The ordered tetrahedral curve is the degree-twelve cover of the
four-point quotient-orbifold deformation curve; its monodromy is
the corresponding pure subgroup of the discrete \(A_4\)
centralizer. The discrete branched-cover normalizer theorem,
or its faithful orbifold Dehn--Nielsen action, makes this
monodromy injective. Hence \(w\) is a nonidentity discrete
tetrahedral braid. B.7 shows that it cannot centralize \(A_5\).
This last assertion concerns the full twelve-point monodromy,
not its four-point images.
\end{proof}

The construction gives one recursive word working for every
four-subset at once. It does not determine the entire solution set
of the paired equations. The equations remain useful necessary
conditions, but eliminating their common kernel is an additional
requirement in route B.4(b).
